\section{Media Queries}

Media queries adalah mekanisme CSS yang memungkinkan penerapan aturan styling berdasarkan kondisi perangkat---terutama \textbf{lebar viewport}. Ini adalah fondasi utama desain responsif: satu halaman HTML yang tampil optimal di desktop, tablet, dan mobile. Sintaks dasar: \code{@media (kondisi) \{ aturan CSS \}}. Kondisi paling umum adalah \code{min-width} dan \code{max-width} \cite{mdnCss}.

\begin{table}[htbp]
\centering
\begin{tabular}{|P{4.5cm}|P{3cm}|P{4.5cm}|}
\hline
\textbf{Media Query} & \textbf{Target} & \textbf{Penggunaan} \\
\hline
\code{max-width: 480px} & Mobile kecil & Font, padding lebih kecil \\
\hline
\code{max-width: 768px} & Tablet portrait & Layout 1--2 kolom \\
\hline
\code{max-width: 1024px} & Tablet landscape & Sidebar collapse \\
\hline
\code{min-width: 1200px} & Desktop besar & Layout penuh, sidebar \\
\hline
\code{prefers-color-scheme} & Tema gelap/terang & Dark mode otomatis \\
\hline
\code{prefers-reduced-motion} & Kurangi animasi & Aksesibilitas \\
\hline
\end{tabular}
\caption{Media queries umum untuk desain responsif}
\end{table}

Selain lebar viewport, media queries juga mendukung fitur lain: \code{orientation} (portrait/landscape), \code{resolution} (untuk layar HiDPI), \code{prefers-color-scheme} (dark/light mode), dan \code{prefers-reduced-motion} (untuk aksesibilitas). Operator logika \code{and}, \code{or} (koma), dan \code{not} memungkinkan kondisi majemuk: \code{@media (min-width: 768px) and (orientation: landscape)} \cite{webdevCss}.

\begin{webcode}[language=CSS, caption={Media queries untuk layout responsif}]
/* Base styles (mobile-first) */
.container { padding: 1rem; }
.grid { display: grid; gap: 1rem; }

/* Tablet: 2 kolom */
@media (min-width: 768px) {
  .grid { grid-template-columns: 1fr 1fr; }
}

/* Desktop: 3 kolom + sidebar */
@media (min-width: 1024px) {
  .grid { grid-template-columns: repeat(3, 1fr); }
  .sidebar { display: block; }
}

/* Dark mode */
@media (prefers-color-scheme: dark) {
  body { background: #1a1a1a; color: #eee; }
}
\end{webcode}

\begin{catatan}
\textbf{Breakpoint Berdasarkan Konten.} Pilih breakpoint berdasarkan kapan layout ``rusak'' (konten terlalu sempit atau terlalu lebar), bukan berdasarkan ukuran perangkat tertentu. Perangkat memiliki ukuran yang sangat beragam; breakpoint berbasis konten lebih tahan masa depan. Uji dengan perlahan menyempitkan browser dan catat di mana layout perlu disesuaikan \cite{cssTricks}.
\end{catatan}

